\subsection{Compatibilidade de roupas e recomendação top-$N$}

Han et al.~\cite{han2017} tratam um look da base Polyvore como uma sequência de peças. Eles treinam uma LSTM bidirecional (Bi-LSTM) que aprende a compatibilidade a partir das imagens e dos textos. Também propõem a prova de completar a lacuna (FITB), em que o modelo escolhe a peça retirada do look entre quatro opções. Vasileva et al.~\cite{vasileva2018} observam que a compatibilidade depende do par de tipos de peça. Por isso, aprendem representações vetoriais (\emph{embeddings}) com subespaços próprios para cada par de tipos. McAuley et al.~\cite{mcauley2015} recomendam produtos por atributos visuais extraídos das imagens e modelam relações de substituição e de complemento entre itens. Os três trabalhos dependem de redes neurais e de imagens. A FashionVerse-MME~\cite{fashionverse} foi criada a partir da Polyvore, o que permite usar a mesma prova FITB.

Deshpande e Karypis~\cite{deshpande2004} propõem a recomendação top-$N$ baseada em itens. A semelhança entre dois itens vem dos usuários que consumiram ambos, e a lista reúne os $N$ itens mais semelhantes aos que o usuário já tem. A avaliação retira um item de cada usuário e verifica se ele volta entre os $N$ primeiros. A métrica Hit@$k$ segue esse protocolo, com o look no lugar do usuário: é a fração de looks em que a peça escondida fica entre as $k$ primeiras.

\subsection{Estruturas usadas}

\textbf{KD-tree.} Bentley~\cite{bentley1975} define a KD-tree como uma árvore binária de busca em várias dimensões. Cada nível divide o espaço por um eixo. Com a mediana em cada nível, a construção custa $O(n \log n)$. Friedman, Bentley e Finkel~\cite{friedman1977} mostram que a busca dos vizinhos mais próximos (kNN) tem custo esperado logarítmico. O motivo é a poda: a busca descarta os lados cuja divisão está mais longe que a pior vizinha já encontrada. No pior caso, a busca visita os $n$ nós. Usamos a KD-tree para achar as peças parecidas com cada peça do look e para montar o grafo kNN dos grupos. Assim, a busca não precisa comparar cada peça com todas as outras.

\textbf{BFS.} A busca em largura visita os vértices por camadas de distância e custa $O(V + E)$. No grafo bipartido de peças e looks, a BFS de 2 passos vai da peça aos looks dela e desses looks às outras peças. Dois passos bastam, porque duas peças a essa distância aparecem juntas em pelo menos um look. O custo é a soma das listas visitadas. O número de vezes que uma peça é alcançada conta os looks em comum, e essa contagem é a base da nota das versões B, C e F.

\textbf{Union-Find.} A estrutura guarda conjuntos disjuntos como árvores. Galler e Fisher~\cite{galler1964} penduram a árvore menor sob a maior, o que limita a altura a $O(\log n)$. A união por posto, usada neste trabalho, segue a mesma ideia, mas usa uma cota superior da altura no lugar do tamanho~\cite{tarjan1975}. Com a compressão de caminho, Tarjan~\cite{tarjan1975} prova que $m$ operações custam $O(m\,\alpha(m, n))$. Aqui, $\alpha$ é a inversa da função de Ackermann e não passa de 4 em tamanhos práticos. No trabalho, o Union-Find guarda os grupos durante o Kruskal e diz, quase em tempo constante, se duas peças já estão no mesmo grupo.

\textbf{Árvore geradora mínima.} O algoritmo de Kruskal~\cite{kruskal1956} ordena as arestas por peso e aceita cada aresta que não forma ciclo. O Union-Find faz esse teste, e o custo $O(E \log E)$ vem da ordenação. O de Prim~\cite{prim1957} cresce uma única árvore a partir de um vértice e, com heap binário, custa $O(E \log V)$. Gower e Ross~\cite{gower1969} mostram que cortar as arestas mais longas da árvore geradora mínima produz os grupos do agrupamento por ligação simples. Zahn~\cite{zahn1971} usa essa árvore para achar grupos. Por isso paramos o Kruskal quando restam $K$ componentes, que já são os grupos. O Prim não permite esse corte no meio, porque cresce uma árvore só. A ligação simples tem um problema conhecido, o efeito corrente: um grupo grande sempre acha mais uma peça perto e absorve as demais.

\textbf{Grafo de condensação e grau.} No grafo de condensação, cada grupo vira um supervértice. Os pontos de combinação passam a ser contados entre grupos de categorias diferentes, e não entre peças. O Tarjan roda nesse grafo. O grau de uma peça no grafo de peças e looks é o número de looks de treino em que ela aparece. A popularidade da versão F é $\log(1 + \text{grau})$.

\textbf{Tarjan.} Tarjan~\cite{tarjan1972} mostra que uma busca em profundidade resolve vários problemas de grafos em tempo linear. Cada vértice recebe a ordem de descoberta e o menor valor que sua subárvore alcança por uma aresta de retorno. Uma aresta é ponte quando a subárvore abaixo dela não alcança nenhum vértice acima. Um vértice é articulação quando sua remoção separa uma subárvore do resto. Hopcroft e Tarjan~\cite{hopcroft1973} publicam esses algoritmos para componentes biconexas. O custo é $O(V + E)$. Usamos o Tarjan, com pilha explícita, para achar os grupos que são a única ligação entre estilos.

\subsection{Montar o look inteiro é NP-difícil}

Montar um look inteiro é escolher uma peça por categoria de modo que a soma das notas entre os pares seja máxima. Esse problema é NP-difícil, e a prova é uma redução do problema da clique, que é NP-completo~\cite{karp1972}. Dado um grafo $G = (V, E)$ e um inteiro $k$, criamos $k$ categorias, cada uma com uma cópia de $V$. As peças $(i, u)$ e $(j, v)$, com $i \neq j$, têm nota 1 se $u$ e $v$ são adjacentes em $G$, e 0 nos outros casos. Como nenhum vértice é adjacente a si mesmo, um look com todas as notas iguais a 1 usa $k$ vértices distintos e adjacentes entre si. Esse conjunto é uma clique de tamanho $k$, e o contrário também vale. A redução cria $k\,|V|$ peças e no máximo $k^2 |V|^2$ notas, em tempo polinomial. Na prática, um look de treino tem em média 4,0 peças. Com as 10 melhores candidatas por categoria, um look de 4 categorias tem $10^4$ combinações, e a enumeração resolve esse número. A avaliação deste trabalho trata da tarefa de completar uma peça por vez.
